% This LaTeX document needs to be compiled with XeLaTeX.
\documentclass[
  10
]{article}
\usepackage[margin=2cm,includehead=true,includefoot=true,centering,]{geometry}
\usepackage{xcolor}
\usepackage{ucharclasses}
\usepackage{hyperref}
\usepackage{amsmath,amssymb}
\usepackage{amsfonts}
\usepackage[version=4]{mhchem}
\usepackage{stmaryrd}
\usepackage{bbold}
\usepackage{polyglossia}
\usepackage{fontspec}
\usepackage{ctex}
\usepackage[export]{adjustbox}
\usepackage{tabularx}
\usepackage{booktabs}

\usepackage{setspace}
\setstretch{1.2}

\makeatletter
\@ifundefined{KOMAClassName}{% if non-KOMA class
  \IfFileExists{parskip.sty}{%
    \usepackage{parskip}
  }{% else
    \setlength{\parindent}{0pt}
    \setlength{\parskip}{6pt plus 2pt minus 1pt}}
}{% if KOMA class
  \KOMAoptions{parskip=half}}
\makeatother
\usepackage{longtable,booktabs,array}
\newcounter{none} % for unnumbered tables
\usepackage{calc} % for calculating minipage widths
% Correct order of tables after \paragraph or \subparagraph
\usepackage{etoolbox}
\makeatletter
\patchcmd\longtable{\par}{\if@noskipsec\mbox{}\fi\par}{}{}
\makeatother
% Allow footnotes in longtable head/foot
\IfFileExists{footnotehyper.sty}{\usepackage{footnotehyper}}{\usepackage{footnote}}
\makesavenoteenv{longtable}
\setlength{\emergencystretch}{3em} % prevent overfull lines
\providecommand{\tightlist}{%
  \setlength{\itemsep}{0pt}\setlength{\parskip}{0pt}}


\hypersetup{colorlinks=true, linkcolor=blue, filecolor=magenta, urlcolor=cyan,}
\urlstyle{same}

\usepackage{colortbl}
\definecolor{table-row-color}{HTML}{999999}
\definecolor{table-rule-color}{HTML}{999999}
%\arrayrulecolor{black!40}
\arrayrulecolor{table-rule-color}     % color of \toprule, \midrule, \bottomrule
\setlength{\aboverulesep}{0pt}
\setlength{\belowrulesep}{0pt}


\setotherlanguages{english}
\IfFontExistsTF{Source Han Serif CN}
{\newfontfamily\chinesefont{Source Han Serif CN}}
{\IfFontExistsTF{Noto Serif CJK SC}
  {\newfontfamily\chinesefont{Noto Serif CJK SC}}
  {\IfFontExistsTF{SimSun}
    {\newfontfamily\chinesefont{SimSun}}
    {\IfFontExistsTF{FangSong}
      {\newfontfamily\chinesefont{FangSong}}
      {\newfontfamily\chinesefont{Arial Unicode MS}}
}}}
\IfFontExistsTF{Times New Roman}
{\newfontfamily\englishfont{Times New Roman}}
{\IfFontExistsTF{Liberation Serif}
  {\newfontfamily\englishfont{Liberation Serif}}
  {\IfFontExistsTF{DejaVu Serif}
    {\newfontfamily\englishfont{DejaVu Serif}}
    {\newfontfamily\englishfont{Arial}}
}}


\author{}
\date{}

\begin{document}


\section{第十八届 (2025 年)}\label{ux7b2cux5341ux516bux5c4a-2025-ux5e74}

\section{北京市大学生物理实验竞赛}\label{ux5317ux4eacux5e02ux5927ux5b66ux751fux7269ux7406ux5b9eux9a8cux7adeux8d5b}

\section{参赛项目说明书}\label{ux53c2ux8d5bux9879ux76eeux8bf4ux660eux4e66}

参赛题目与类别：（请在□处画√）

\subsection{命题类}\label{ux547dux9898ux7c7b}

□题目1：微小位移测量

□题目2：探究电磁感应现象中的能量转换

□题目3：弱压力测量

□题目4：晶体双折射

\subsection{制作类}\label{ux5236ux4f5cux7c7b}

□题目5：实验仪器的制作、改进（□定量类 □演示类）

□题目6：前沿物理（□前沿文献类 □自研文献类）

□题目7：AI+物理实验

\subsection{信息资源类}\label{ux4fe1ux606fux8d44ux6e90ux7c7b}

□题目8：大学物理教学微视频

□题目9：物理教学资源开发（□媒体资源类 □仿真程序类）

北京市大学生物理实验竞赛组委会制

2025年6月

{\def\LTcaptype{none} % do not increment counter
\begin{longtable}[]{@{}|l|l|@{}}
\toprule\noalign{}
\endhead
\bottomrule\noalign{}
\endlastfoot
\hline
参赛题目类型(题目5、6、9填写) & \\
\hline
题目名称 & \\
\hline
\multicolumn{2}{@{}l@{}}{%
(以下内容可加页,篇幅3-5页)设计原理与方法:实验仪器与装置:数据测量与分析:结论:项目创新与特色:制作成本(明细)} \\
\hline
\end{longtable}
}

注意:

竞赛采用匿名评审，说明书里不允许出现任何参赛队和参赛学校信息。

邮箱: wlsyzx@bjtu.edu.cn 联系电话: 51683629, 13811861683 贺敬文

通信地址：北京交通大学物理国家级实验教学示范中心 邮编:100044

\end{document}
